\subsection{The actual option and the hedged price need separate calculations}
\label{actual-aligned-option-components}

The time-value comparison also has to fit the risk being hedged. An actual
option may have a different exercise price from the price level the company
has designated. Its payoff then differs from the one-sided change in the
hedged purchase. IFRS~9.B6.5.5 permits a hypothetical derivative to measure
that item but prohibits importing features that belong only to the actual
instrument. B6.5.32 similarly measures aligned time value using an option
whose critical terms match the item
\citep[IFRS~9.B6.5.5, B6.5.32--33]{eu20231803}. The following complete
transaction uses the retained historical EU compilation; it is a distinct
conditional example from the preceding AASB-profile illustrations.

Consider a one-unit inventory purchase with a designated one-sided price risk
above 110. The company buys a call with strike 100. At designation, spot is 90
and both calls have zero intrinsic value. The actual option costs 6, entirely
time value; the hypothetical call with strike 110 has aligned time value 4.
These component values and the qualifying designation are supplied. The
company pays only the actual premium. The hypothetical option is a calculation
of the hedged risk, not another acquired asset.

At a later spot of 115, actual intrinsic value is \(115-100=15\), while the
hedged item's aligned intrinsic value is \(115-110=5\). Suppose actual and
aligned time values are both 3. The actual option is worth 18 and the
hypothetical option is worth 8. The effective intrinsic gain is 5; the
remaining actual intrinsic gain of 10 enters profit. Because initial actual
time value exceeded aligned time value, B6.5.33(a) gives a cost reserve change
of \(3-4=-1\) and time-value mismatch profit of
\((3-6)-(-1)=-2\). Thus total reserve is 4 and profit is 8, together
accounting for the actual option's gain of \(18-6=12\).

The distinction can be expressed directly from the contractual price levels.
Let \(x\) be the observed price, \(q\) the common quantity, \(K_A\) the
actual strike and \(K_H\) the hedged price level. The two intrinsic values are
\[
 I_A=q\max(x-K_A,0),\qquad I_H=q\max(x-K_H,0).
\]
With zero initial intrinsic values, these are also the cumulative intrinsic
gains. In the supplied qualifying relationship, the intrinsic reserve is
\(R_I=\min(I_A,I_H)\), and intrinsic mismatch profit is \(I_A-R_I\).
Writing actual and aligned time values as \(T_A,T_H\), with designation
values \(T_{A0},T_{H0}\), the admitted cumulative-loss cases give
\[
 R_T=
 \begin{cases}
 T_H-T_{H0},&T_{A0}>T_{H0},\\
 -\min(T_{A0}-T_A,T_{H0}-T_H),&T_{A0}<T_{H0}.
 \end{cases}
 \qquad
 \Pi=(I_A-R_I)+(T_A-T_{A0}-R_T).
\]
Here \(R_T\) is the separate cost reserve. Both time values remain between
zero and their own initial values in this profile; the second branch compares
cumulative losses. It is not a signed-value rule for every collar or opposed
change. Matching initial values are admitted here only with fully matched
terms and matching subsequent component values. The intrinsic payoff identity
does not price either time value.

At expiry spot 130, actual intrinsic value is 30 and aligned intrinsic value
is 20, while both time values are zero. The intrinsic reserve is 20, the cost
reserve is \(-4\), and cumulative profit remains \(10-2=8\). Start with
cash and capital 200: pay premium 6, pay the inventory supplier 130, and receive
actual option settlement 30. Cash closes at 94. The two basis transfers give
\[
 \text{Inventory cost}=130-20-(-4)=114,\qquad
 94+114=200+8.
\]
The intrinsic reserve reduces the inventory cost, while the negative cost
reserve adds the related protection cost. Both transfers are separate from
cash under IFRS~9.6.5.11(d)(i) and 6.5.15(b)(i)
\citep[IFRS~9.6.5.11, 6.5.15]{eu20231803}. If actual intrinsic value 30
were copied into the hedged-item calculation, inventory would become 104 and
profit \(-2\). Cash 94 plus inventory 104 would still equal capital 200 less
the loss 2. The book would balance while measuring the wrong hedged risk.

For the lower-actual-time-value branch, reverse the strikes: actual strike
110, hedged price 100, and initial actual/aligned time values 2 and 3. At spot
105, actual intrinsic is zero and aligned intrinsic is 5. With time values
1 and \(5/2\), the smaller cumulative loss is \(1/2\); the cost reserve
and mismatch profit are each \(-1/2\). At spot 115, time values \(1/2\)
and 1 give cumulative losses \(3/2\) and 2. The cost reserve is now
\(-3/2\) and cumulative time mismatch profit is zero. The intervening
profit gain of \(1/2\) reverses the earlier mismatch. Applying a new
lower-of comparison only to each period would leave a reserve of \(-1\)
and miss that reversal.

At expiry spot 120, actual/aligned intrinsic values are 10 and 20, and time
values are zero. The intrinsic reserve is 10, the cost reserve is \(-2\),
and cumulative profit is zero. From capital 200, cash closes at
\(200-2-120+10=88\) and inventory at \(120-10+2=112\). The full
premium, supplier payment and actual settlement are again visible; the
hypothetical payoff of 20 creates neither an asset nor cash.

The read-only \texttt{AlignedOptionComponents/1} report binds each observation
to both contract identities and checks its intrinsic/time decomposition. Six
paths include matched terms, out-of-the-money expiry and doubled quantities.
Unknown component values, copied critical terms, omitted premiums, mixed rule
profiles and unsupported cumulative changes produce no partial report. Actual
pricing, qualifying findings, general signed components and subsequent company
or sector integration retain their separate requirements.
